\subsection{Random Variables}

\hrulefill

\vspace{5mm}

In many random experiments, we are not interested in the specific outcome of an experiment, but in some characteristic of
that outcome. If you play the lottery , for example, you would like the numbers chosen in the weekly drawing to match
as many of the number you marked on your ticket as possible. You do not particularly care which of your numbers are
matched, but rather how many matches you have.

\paragraph*{Definition}
For a given sample space of some experiment, a \textbf{random variable} is any rule that associates a number with 
each outcome in S. In mathematical language, a random variable is a function whose domain is the sample space S and whose
range is the set of real numbers. We usually denote random variables with capital letters such as X, Y, and Z.

\begin{align*}
    X: S &\to \mathbb{R} \\
    \exists s \in S \ni X(s) &= x
\end{align*}

\paragraph*{Random Variable Examples}
\begin{itemize}
    \item Sum of two dice
    \item Number of defects on a randomly selected piece of fabric.
    \item SAT score of a randomly selected student.
    \item Number of heads in 10 flips of a coin.
    \item The amount a company has to pay to an employee in pension.
\end{itemize}

\paragraph*{Types of Random Variables}
\begin{itemize}
    \item A \textbf{discrete random variable} is a random variable that can take on a finite or countably infinite
    number of possible values.
    \item A \textbf{continuous random variable} is a random variable that can take on infinitely many values corresponding
    to an interval on the real line.
\end{itemize}

\subsubsection*{Probability Distributions}
\paragraph*{Definition}
To describe a discrete random variable X, you need two pieces of information:
\begin{itemize}
    \item What values X can assume.
    \item The probabilities associated with each of these values.
\end{itemize}

The \textbf{probability distribution} or \textbf{probability mass function} (pmf) of a discrete random variable X
is a formula, table, or graph that gives both the possible values of X and the probabilities associated with each value 
of X. The probability distribution or pmf of a discrete random variable is defined for every number x by:

\[
    p(x) = P(X = x) P(\forall s \in S, X(s) = x)
\]

\paragraph*{Example:}
Let X be the random variable that gives the sum of 2d6.
\begin{table}[h!]
    \centering
    \begin{tabular}{c|c|c|c|c|c|c|c|c|c|c|c|c}
        x & 2 & 3 & 4 & 5 & 6 & 7 & 8 & 9 & 10 & 11 & 12 \\
        \hline
        p(x) & $\frac{1}{36}$ & $\frac{2}{36}$ & $\frac{3}{36}$ & $\frac{4}{36}$ & $\frac{5}{36}$ & $\frac{6}{36}$ & $\frac{5}{36}$ & $\frac{4}{36}$ & $\frac{3}{36}$ & $\frac{2}{36}$ & $\frac{1}{36}$
    \end{tabular}
\end{table}

\pagebreak

\subsubsection*{Bernoulli Random Variables}
\paragraph*{Definition}
Any random variable whose only possible values are 0 and 1 is called a \textbf{Bernoulli random variable}.
The PMF looks like this:
\begin{table}[h!]
    \centering
    \begin{tabular}{c|c|c}
        x & 0 & 1 \\
        \hline
        p(x) & $1-p$ & $p$
    \end{tabular}
\end{table}

\paragraph*{Example:}
A box contains one red, two orange, and five black balls. You choose two balls at random. For each red ball you win \$5,
for each orange ball you win \$1, and for each black ball you win nothing. Let X be the random variable that gives your 
total winnings. Find the probability distribution of X and the probability that you will win at most \$3.

\begin{table}[h!]
    \centering
    \tabular{|c|c|c|}
        \hline
        x & p(x) & Explanation \\
        \hline
        0 & $\frac{10}{28}$ & BB \\
        \hline
        1 & $\frac{10}{28}$ & BO, OB \\
        \hline
        2 & $\frac{1}{28}$ & OO \\
        \hline
        5 & $\frac{5}{28}$ & RB, BR \\
        \hline
        6 & $\frac{2}{28}$ & RO, OR \\
        \hline
    \endtabular
\end{table}

\[
    P(X \leq 3) = p(0) + p(1) + p(2) = \frac{10}{28} + \frac{10}{28} + \frac{1}{28} = \frac{21}{28} = \frac{3}{4}
\]

\pagebreak

\subsubsection*{Cumulative Distribution Function}
\paragraph*{Definition}
The \textbf{cumulative distribution function} (CDF) $F(x)$ of a discrete RV $X$ with PMF $p(x)$ is defined for \textbf{every number x} by:
\[
    F(x) = P(X \leq x) = \sum_{i \leq x} p(i)
\]

\paragraph*{Example:}
Take the previous example with the balls. Some examples of using CDF here are:
\begin{align*}
    F(3) &= P(X \leq 3) = p(0) + p(1) + p(2) = \frac{21}{28} = \frac{3}{4} \\
    F(5) &= P(X \leq 5) = p(0) + p(1) + p(2) + p(5) = \frac{26}{28} = \frac{13}{14} \\
    F(-10) &= P(X \leq -10) = 0 \\
    F(10) &= P(X \leq 10) = 1
\end{align*}

We can find the CDF:
\begin{align*}
    F(x) = 
    \begin{cases}
        0 & x < 0 \\
        \frac{10}{28} & 0 \leq x < 1 \\
        \frac{20}{28} & 1 \leq x < 2 \\
        \frac{21}{28} & 2 \leq x < 5 \\
        \frac{26}{28} & 5 \leq x < 6 \\
        1 & x \geq 6
    \end{cases}
\end{align*}

